\documentclass[a4paper,11pt]{article}

\usepackage{revy}
\usepackage[utf8]{inputenc}
\usepackage[T1]{fontenc}
\usepackage[danish]{babel}


\revyname{MatematikRevyen}
\revyyear{2026}
% HUSK AT OPDATERE VERSIONSNUMMER
\version{0.1}
\eta{$1$ minutter}
\status{Færdig}

\title{Kim Possible}
\author{Victoria 19'}

\begin{document}
\maketitle

\begin{roles}
\role{K}[Skuespiller/Sanger] Kim Possible/instruktor 
\role{S1}[Skuespiller] Studerende
\role{S2}[Skuespiller] Studerende/statist
\role{S3}[Skuespiller] Studerende/statist
\end{roles}

\begin{props}
\prop{Rekvisit} Tavle og kridt
\prop{Rekvisit} Borde og stole til de studerende
\end{props}


\begin{sketch}
\scene{Første øvelsestime i An$0$. De 3 elever sidder ved deres borde og har bogen fremme. K stiller sig hen til tavlen.}

\says{K} Godmorgen allesammen! Jeg er jeres instruktor i Analyse $0$, så det er mig, som kommer til at rette jeres afleveringer og gennemgå øvelserne for jer hver uge. Hvis I har nogle spørgsmål, eller hvis der er noget I ikke forstår, så skal I ikke tøve med at række hånden op. 

\scene{S1 rækker hånden op}

\says{K} Ja, hvad vil du spørge om??

\says{S1} Hvad hvis vi har spørgsmål uden for øvelestimen? 

\says{K} Jamen så kan I bare 

\scene{Bandet spiller omkv fra Kim Possible (fra 0:26 på Spotify's 'Mail mig, ring mig - Fra "Kim Possible"' af Saseline Sørensen}

\says{K} Mail mig, ring mig, det gør ingenting nej \\
Sms er også helt okay\\
Så hjælper jeg gerne baby\\
Mail mig, ring mig, det gør ingenting nej\\ \\
\scene{K flipper tavlen, og bagpå står KU-mail og et tlf-nr (evt. revyens mobilepay nr). De studerende skriver det ned}
\says{K} Mail mig, ring mig, det gør ingenting nej \\
Ligemeget hvor, ligemeget hvornår\\
Er du sikker på at hjælpen du vil få\\
Problemer? Jeg svarer\\
Det hele, vi klarer\\
Du ved at du altid kan nå\\
Kim Possible\\

\says{K} Har du problemer?

\says{K} Mail mig, ring mig, det gør ingenting nej 


\scene{Tæppe}

\end{sketch}
\end{document}
